\documentclass[11pt]{article}
\newcommand{\ListaTitulo}{Gabarito --- Lista 06}
% Preâmbulo compartilhado — MATD48, Listas de Exercícios 2026
% Usado por Lista{NN}.tex e Gabarito{NN}.tex via % Preâmbulo compartilhado — MATD48, Listas de Exercícios 2026
% Usado por Lista{NN}.tex e Gabarito{NN}.tex via % Preâmbulo compartilhado — MATD48, Listas de Exercícios 2026
% Usado por Lista{NN}.tex e Gabarito{NN}.tex via \input{preamble}

\usepackage[utf8]{inputenc}
\usepackage[T1]{fontenc}
\usepackage[brazil]{babel}
\usepackage{fullpage}
\usepackage{amsmath,amssymb}
\usepackage{enumitem}
\usepackage{xcolor}
\usepackage{fancyhdr}
\usepackage{titlesec}
\usepackage{listings}
\usepackage{hyperref}
\usepackage{booktabs}
\usepackage{graphicx}

\definecolor{matd48azul}{HTML}{0B4F6C}
\definecolor{matd48verde}{HTML}{2E8B57}

\titleformat{\section}{\normalfont\Large\bfseries\color{matd48azul}}{\thesection}{1em}{}
\titleformat{\subsection}{\normalfont\large\bfseries\color{matd48azul}}{\thesubsection}{1em}{}

\providecommand{\ListaTitulo}{Lista de Exercícios}

\setlength{\headheight}{14pt}
\pagestyle{fancy}
\fancyhf{}
\lhead{\small MATD48 --- Planejamento de Experimentos A}
\rhead{\small \ListaTitulo}
\cfoot{\thepage}
\renewcommand{\headrulewidth}{0.4pt}

\lstset{
  basicstyle=\ttfamily\small,
  breaklines=true,
  frame=single,
  columns=fullflexible,
  backgroundcolor=\color{gray!8},
  commentstyle=\color{matd48verde},
  keywordstyle=\color{matd48azul}\bfseries,
  language=R,
  extendedchars=true,
  literate=%
    {á}{{\'a}}1 {à}{{\`a}}1 {â}{{\^a}}1 {ã}{{\~a}}1
    {é}{{\'e}}1 {ê}{{\^e}}1 {í}{{\'i}}1 {ó}{{\'o}}1
    {ô}{{\^o}}1 {õ}{{\~o}}1 {ú}{{\'u}}1 {ç}{{\c c}}1
    {Á}{{\'A}}1 {À}{{\`A}}1 {Â}{{\^A}}1 {Ã}{{\~A}}1
    {É}{{\'E}}1 {Ê}{{\^E}}1 {Í}{{\'I}}1 {Ó}{{\'O}}1
    {Ô}{{\^O}}1 {Õ}{{\~O}}1 {Ú}{{\'U}}1 {Ç}{{\c C}}1
}

\hypersetup{colorlinks=true, linkcolor=matd48azul, urlcolor=matd48azul, citecolor=matd48azul}

\newcommand{\cabecalho}[2]{%
  \begin{center}
    {\Large\bfseries\color{matd48azul} #1}\\[0.3em]
    {\large #2}\\[0.3em]
    \rule{\linewidth}{0.6pt}
  \end{center}
  \vspace{0.5em}
}

\newenvironment{questao}[1]{\vspace{1em}\noindent\textbf{Questão #1.}\hspace{0.4em}}{\par}
\newenvironment{solucao}{\vspace{0.3em}\noindent\textit{Solução.}\hspace{0.4em}\color{black!85}}{\par\vspace{0.5em}}


\usepackage[utf8]{inputenc}
\usepackage[T1]{fontenc}
\usepackage[brazil]{babel}
\usepackage{fullpage}
\usepackage{amsmath,amssymb}
\usepackage{enumitem}
\usepackage{xcolor}
\usepackage{fancyhdr}
\usepackage{titlesec}
\usepackage{listings}
\usepackage{hyperref}
\usepackage{booktabs}
\usepackage{graphicx}

\definecolor{matd48azul}{HTML}{0B4F6C}
\definecolor{matd48verde}{HTML}{2E8B57}

\titleformat{\section}{\normalfont\Large\bfseries\color{matd48azul}}{\thesection}{1em}{}
\titleformat{\subsection}{\normalfont\large\bfseries\color{matd48azul}}{\thesubsection}{1em}{}

\providecommand{\ListaTitulo}{Lista de Exercícios}

\setlength{\headheight}{14pt}
\pagestyle{fancy}
\fancyhf{}
\lhead{\small MATD48 --- Planejamento de Experimentos A}
\rhead{\small \ListaTitulo}
\cfoot{\thepage}
\renewcommand{\headrulewidth}{0.4pt}

\lstset{
  basicstyle=\ttfamily\small,
  breaklines=true,
  frame=single,
  columns=fullflexible,
  backgroundcolor=\color{gray!8},
  commentstyle=\color{matd48verde},
  keywordstyle=\color{matd48azul}\bfseries,
  language=R,
  extendedchars=true,
  literate=%
    {á}{{\'a}}1 {à}{{\`a}}1 {â}{{\^a}}1 {ã}{{\~a}}1
    {é}{{\'e}}1 {ê}{{\^e}}1 {í}{{\'i}}1 {ó}{{\'o}}1
    {ô}{{\^o}}1 {õ}{{\~o}}1 {ú}{{\'u}}1 {ç}{{\c c}}1
    {Á}{{\'A}}1 {À}{{\`A}}1 {Â}{{\^A}}1 {Ã}{{\~A}}1
    {É}{{\'E}}1 {Ê}{{\^E}}1 {Í}{{\'I}}1 {Ó}{{\'O}}1
    {Ô}{{\^O}}1 {Õ}{{\~O}}1 {Ú}{{\'U}}1 {Ç}{{\c C}}1
}

\hypersetup{colorlinks=true, linkcolor=matd48azul, urlcolor=matd48azul, citecolor=matd48azul}

\newcommand{\cabecalho}[2]{%
  \begin{center}
    {\Large\bfseries\color{matd48azul} #1}\\[0.3em]
    {\large #2}\\[0.3em]
    \rule{\linewidth}{0.6pt}
  \end{center}
  \vspace{0.5em}
}

\newenvironment{questao}[1]{\vspace{1em}\noindent\textbf{Questão #1.}\hspace{0.4em}}{\par}
\newenvironment{solucao}{\vspace{0.3em}\noindent\textit{Solução.}\hspace{0.4em}\color{black!85}}{\par\vspace{0.5em}}


\usepackage[utf8]{inputenc}
\usepackage[T1]{fontenc}
\usepackage[brazil]{babel}
\usepackage{fullpage}
\usepackage{amsmath,amssymb}
\usepackage{enumitem}
\usepackage{xcolor}
\usepackage{fancyhdr}
\usepackage{titlesec}
\usepackage{listings}
\usepackage{hyperref}
\usepackage{booktabs}
\usepackage{graphicx}

\definecolor{matd48azul}{HTML}{0B4F6C}
\definecolor{matd48verde}{HTML}{2E8B57}

\titleformat{\section}{\normalfont\Large\bfseries\color{matd48azul}}{\thesection}{1em}{}
\titleformat{\subsection}{\normalfont\large\bfseries\color{matd48azul}}{\thesubsection}{1em}{}

\providecommand{\ListaTitulo}{Lista de Exercícios}

\setlength{\headheight}{14pt}
\pagestyle{fancy}
\fancyhf{}
\lhead{\small MATD48 --- Planejamento de Experimentos A}
\rhead{\small \ListaTitulo}
\cfoot{\thepage}
\renewcommand{\headrulewidth}{0.4pt}

\lstset{
  basicstyle=\ttfamily\small,
  breaklines=true,
  frame=single,
  columns=fullflexible,
  backgroundcolor=\color{gray!8},
  commentstyle=\color{matd48verde},
  keywordstyle=\color{matd48azul}\bfseries,
  language=R,
  extendedchars=true,
  literate=%
    {á}{{\'a}}1 {à}{{\`a}}1 {â}{{\^a}}1 {ã}{{\~a}}1
    {é}{{\'e}}1 {ê}{{\^e}}1 {í}{{\'i}}1 {ó}{{\'o}}1
    {ô}{{\^o}}1 {õ}{{\~o}}1 {ú}{{\'u}}1 {ç}{{\c c}}1
    {Á}{{\'A}}1 {À}{{\`A}}1 {Â}{{\^A}}1 {Ã}{{\~A}}1
    {É}{{\'E}}1 {Ê}{{\^E}}1 {Í}{{\'I}}1 {Ó}{{\'O}}1
    {Ô}{{\^O}}1 {Õ}{{\~O}}1 {Ú}{{\'U}}1 {Ç}{{\c C}}1
}

\hypersetup{colorlinks=true, linkcolor=matd48azul, urlcolor=matd48azul, citecolor=matd48azul}

\newcommand{\cabecalho}[2]{%
  \begin{center}
    {\Large\bfseries\color{matd48azul} #1}\\[0.3em]
    {\large #2}\\[0.3em]
    \rule{\linewidth}{0.6pt}
  \end{center}
  \vspace{0.5em}
}

\newenvironment{questao}[1]{\vspace{1em}\noindent\textbf{Questão #1.}\hspace{0.4em}}{\par}
\newenvironment{solucao}{\vspace{0.3em}\noindent\textit{Solução.}\hspace{0.4em}\color{black!85}}{\par\vspace{0.5em}}


\begin{document}

\cabecalho{Gabarito --- Lista de Exercícios 06}{Verificação de pressupostos e transformações no DCA (Capítulo 3 / Aula 06)}

\begin{questao}{1} \textbf{(Psicologia --- interpretando um Q-Q plot)}
\end{questao}
\begin{solucao}
\begin{enumerate}[label=(\alph*)]
  \item Esse padrão --- cauda superior mais pesada que a normal, cauda inferior normal --- é a
    assinatura de \textbf{assimetria à direita}: alguns tempos de reação bem mais lentos que o
    esperado sob normalidade, sem um efeito espelhado do lado dos tempos rápidos. É exatamente o
    formato típico de dados de tempo de reação: há um limite físico inferior (não se reage
    instantaneamente), mas nenhum limite superior firme, produzindo uma cauda longa à direita.
  \item Não necessariamente contradiz: com $N=45$, o teste de Shapiro-Wilk tem poder moderado, e
    pode não rejeitar mesmo havendo assimetria real e visível no Q-Q plot --- ausência de
    significância não é evidência de que o pressuposto vale (Aula 06, Teoria). A leitura correta
    é priorizar a inspeção gráfica quando ela mostra um padrão sistemático e interpretável, usando
    o teste formal como informação complementar, não como árbitro único, especialmente em amostras
    pequenas/moderadas.
  \item A transformação \textbf{logarítmica} ($\log Y$), primeira candidata para dados
    assimétricos à direita, estritamente positivos, com variância que tende a crescer com a
    média --- o padrão clássico de tempos de reação (ver também Questão 3).
\end{enumerate}
\end{solucao}

\begin{questao}{2} \textbf{(Psicologia --- homogeneidade de variância)}
\end{questao}
\begin{solucao}
\begin{enumerate}[label=(\alph*)]
  \item Sim: $p=0{,}0012 < 0{,}05$, rejeita-se $H_0$ de variâncias iguais. Isso compromete a
    validade dos ICs e do teste $F$ da Aula 04, que assumem explicitamente um único $\sigma^2$
    comum a todos os grupos (usado para combinar/\emph{pool} o $MS_E$); com heterocedasticidade
    real, o $MS_E$ combinado não estima corretamente a variância de nenhum grupo específico, e as
    taxas de erro Tipo I/II nominais deixam de valer.
  \item $s^2_{36h}/s^2_{0h} = 81^2/30^2 = 6561/900 \approx \textbf{7{,}29}$ --- uma razão bem
    acima do que se costuma considerar tolerável (regra prática comum: razão $<4$). O padrão
    (variância crescendo junto com a severidade/o nível do fator, não constante) é consistente com
    \textbf{erro multiplicativo}, não aditivo: sob erro aditivo puro a variância seria a mesma em
    todos os grupos, independentemente do nível de privação de sono.
  \item A transformação \textbf{log} é a candidata natural: como visto na Questão 3, sob erro
    multiplicativo a variância na escala original é proporcional ao quadrado da média
    ($\mathrm{Var}(Y)\propto\mu^2$), exatamente o padrão observado aqui (grupos com tempo médio
    maior --- mais privação de sono --- têm desvio-padrão maior); a escala log converte esse erro
    multiplicativo em aditivo de variância constante.
\end{enumerate}
\end{solucao}

\begin{questao}{3} \textbf{(Dedução --- por que o log estabiliza a variância sob erro multiplicativo)}
\end{questao}
\begin{solucao}
\begin{enumerate}[label=(\alph*)]
  \item Aplicando $\log$ a ambos os lados de $Y_{ij}=\mu_i\exp(\varepsilon_{ij})$:
    $\log Y_{ij} = \log\mu_i + \log\exp(\varepsilon_{ij}) = \log\mu_i + \varepsilon_{ij}$.
    Definindo $\mu_i^{\log} := \log\mu_i$, isso é literalmente $\log Y_{ij} = \mu_i^{\log} +
    \varepsilon_{ij}$, com $\varepsilon_{ij}\overset{iid}{\sim}N(0,\sigma^2)$ --- o modelo aditivo
    padrão do DCA (Aula 04), com variância $\sigma^2$ igual para todo $i$ (não depende do índice
    de grupo em nenhum lugar da expressão).
  \item Como $\mu_i$ é uma constante (não aleatória) e $Y_{ij}=\mu_i\cdot W_{ij}$ com
    $W_{ij}=\exp(\varepsilon_{ij})$, $\mathrm{Var}(Y_{ij}) = \mu_i^2\,\mathrm{Var}(W_{ij})$.
    Usando o fato fornecido, $\mathrm{Var}(W_{ij}) = (e^{\sigma^2}-1)e^{\sigma^2} =: C$, uma
    constante que depende só de $\sigma^2$ (igual para todo $i$, pois $\varepsilon_{ij}$ tem a
    mesma distribuição $N(0,\sigma^2)$ em todos os grupos). Logo $\mathrm{Var}(Y_{ij}) = C\,\mu_i^2$.
  \item Pelo item (b), grupos com média $\mu_i$ maior têm, na escala original, variância
    proporcional a $\mu_i^2$ --- maior média implica maior variância, exatamente o padrão
    observado na Questão 2. Pelo item (a), a transformação log remove essa dependência: na escala
    $\log Y$, todo grupo tem a mesma variância $\sigma^2$, restaurando a homocedasticidade exigida
    pela ANOVA.
\end{enumerate}
\end{solucao}

\begin{questao}{4} \textbf{(Ciência de dados --- interpretando Box-Cox)}
\end{questao}
\begin{solucao}
\begin{enumerate}[label=(\alph*)]
  \item O intervalo $(-1{,}1,\ 0{,}3)$ \textbf{inclui} $\lambda=0$ (log é estatisticamente
    compatível com os dados), mas \textbf{não inclui} $\lambda=1$ (não transformar é rejeitado
    pelo procedimento --- alguma transformação é recomendável).
  \item A transformação \textbf{log} ($\lambda=0$): está dentro do intervalo de verossimilhança
    (portanto não é uma escolha pior, em termos estatísticos, que o $\hat\lambda=-0{,}4$ que
    maximiza a verossimilhança), e é muito mais interpretável --- diferenças na escala log
    correspondem a razões multiplicativas na escala original, uma leitura natural para tempos de
    carregamento, enquanto $\lambda=-0{,}4$ não tem interpretação prática direta.
  \item Como nenhum valor "redondo" está no intervalo, as opções são: (i) usar o próprio
    $\hat\lambda=-2{,}15$ (ponto médio do intervalo, ou o valor que de fato maximiza a
    verossimilhança), aceitando a perda de interpretabilidade direta; (ii) usar o valor "redondo"
    mais próximo permitido pelo intervalo, aqui $\lambda=-2$ (transformação $1/Y^2$), como
    compromisso entre ajuste estatístico e interpretabilidade; (iii) investigar se a assimetria
    residual reflete um problema mais estrutural nos dados (outliers, mistura de subpopulações)
    que nenhuma transformação de potência resolveria bem, considerando alternativas não
    paramétricas (Aula 07 em diante).
\end{enumerate}
\end{solucao}

\begin{questao}{5} \textbf{(Integradora --- decidindo o que fazer)}
\end{questao}
\begin{solucao}
\begin{enumerate}[label=(\alph*)]
  \item Primeiro, inspeção gráfica dos resíduos: Q-Q plot (normalidade) e resíduos vs. valores
    ajustados (padrão de variância, presença de tendência). Em seguida, testes formais de apoio
    --- Shapiro-Wilk (normalidade) e Levene (homocedasticidade) --- lembrando que eles confirmam,
    não substituem, a leitura gráfica. Só depois desse diagnóstico decide-se por transformar os
    dados ou não.
  \item A transformação \textbf{raiz quadrada} ($\sqrt{Y}$, ou $\sqrt{Y+1/2}$ se houver zeros) é a
    candidata natural para dados de \textbf{contagem}, nos quais a variância tende a crescer
    proporcionalmente à \emph{média} (não ao seu quadrado, como no caso multiplicativo/log da
    Questão 3) --- um padrão consistente com a distribuição de Poisson, frequentemente razoável
    para contagens de eventos raros como erros em uma tarefa.
  \item Se, mesmo após transformar, os resíduos continuarem severamente não normais, a
    pesquisadora deveria considerar uma alternativa \textbf{não paramétrica}, como o teste de
    Kruskal-Wallis. Conceitualmente, isso muda o que está sendo comparado: em vez de médias
    populacionais sob um modelo distribucional específico, compara-se a distribuição das
    \textbf{ordens (postos)} dos erros entre as condições de iluminação --- uma pergunta sobre
    posição relativa das distribuições, válida sob suposições muito mais fracas, ao custo de
    alguma perda de poder quando o modelo normal (ou uma transformação dele) de fato seria
    razoável.
\end{enumerate}
\end{solucao}

\end{document}
